\section{Introduction}

Interface Definition Languages (IDLs) such as Google Protocol Buffers (Protobuf) \cite{protobuf_spec}, FlatBuffers \cite{flatbuffers}, and Cap'n Proto \cite{capnproto} have become foundational to distributed systems, microservice architectures, and high-performance computing. Protocol Buffers provide cross-language interoperability, compact wire encodings, and robust backward- and forward-compatibility guarantees.

Historically, Protocol Buffers rely on an external code generator, \texttt{protoc}. When compiling a C++ project that utilizes Protobuf schemas, the build system must execute \texttt{protoc} to generate intermediate \texttt{.pb.h} header files and \texttt{.pb.cc} implementation files. While functional, this workflow introduces substantial engineering friction:

\begin{itemize}
    \item \textbf{Build System Complexity:} Incorporating external code generators requires intricate build rules (such as CMake \texttt{add\_custom\_command} or Bazel rules), tracking dependency directed acyclic graphs (DAGs), and synchronizing generator toolchains across development and production environments.
    \item \textbf{Version Skew and Toolchain Synchronization:} Mismatches between the runtime \texttt{libprotobuf} library and the build-time \texttt{protoc} compiler binary frequently produce subtle ABI incompatibilities or compilation failures.
    \item \textbf{Codebase Bloat and Intermediate Artifacts:} For large enterprise systems with hundreds of messages, generated header files span hundreds of thousands of lines of code, placing significant memory pressure on compiler front-ends and slowing incremental builds.
    \item \textbf{Loss of Direct Language Integration:} Developers must maintain external build targets rather than expressing message schemas as native, compile-time entities within the C++ type system.
\end{itemize}

With the upcoming ISO C++26 standard, the language introduces two transformative features: P1967 \texttt{\#embed} \cite{p1967} for zero-overhead compile-time file inclusion, and P2996 Static Reflection \cite{p2996} with splicing operators for programmatic type synthesis. 

In this paper, we present \texttt{ctpb} (Compile-Time Protobuf) \cite{ctpb2026}, a header-only modern C++26 library that eliminates the external \texttt{protoc} code generator entirely. \texttt{ctpb} embeds \texttt{.proto} schema files directly into the translation unit at compile time, executes a fully \texttt{constexpr} recursive-descent parser inside the compiler front-end, and synthesizes type-safe C++ structs, accessors, and wire-format serialization routines through static reflection.

The primary contributions of this work are:
\begin{enumerate}
    \item We design a three-stage compile-time architecture that transforms raw IDL text into first-class C++ types entirely within the compiler front-end without intermediate disk I/O.
    \item We implement a zero-allocation, \texttt{constexpr} recursive-descent proto3 parser capable of validating schemas, parsing nested messages, enums, maps, oneofs, and reserved ranges at compile time.
    \item We demonstrate the application of C++26 P2996 static reflection to programmatically inject member fields, wire visitors, and equality operators.
    \item We evaluate \texttt{ctpb}'s compile-time and runtime performance, showing schema parsing latencies under $11.4\,\mu\text{s}$ and wire serialization throughput exceeding $300\,\text{MB/s}$.
\end{enumerate}
